\documentclass[10pt,a4paper]{amsart}
\usepackage[margin=19mm]{geometry}
\usepackage{amsmath,amssymb,mathtools}
\usepackage[scr=rsfs,cal=euler]{mathalfa}
\usepackage{xfrac}
\usepackage[unicode,hidelinks,bookmarks=false]{hyperref}
\newcommand{\ie}{i.e.\ }
\newcommand{\cf}{cf.\ }
\newcommand{\resp}{respectively\ }
\newcommand{\Subsection}{\subsection}
\providecommand{\nobreakdash}{\nobreak\hbox{-}}
\setlength{\emergencystretch}{2em}
\multlinegap0pt
\usepackage{amssymb}
\usepackage{dsfont}
\newtheorem{lemma}{Lemma}
\newcommand{\cbnorm}[1]{\left\| #1 \right\|_{\mathrm{cb}}}
\newcommand{\complexes}{\mathbb{C}}
\newcommand{\id}[1]{\mathrm{id}_{#1}}
\newcommand{\matr}[1]{\mathbb{M}_{#1}(\complexes)}
\newcommand{\matra}[2]{\mathbb{M}_{#1}(#2)}
\newcommand{\naturals}{\mathbb{N}}
\newcommand{\spec}[1]{\operatorname{\sigma}{\left(#1\right)}}
\newcommand{\unit}{\mathds{1}}
\title{On a generalization of decomposable maps on C*-algebras}
\author{Krzysztof Szczygielski}
\date{Source: arXiv:2602.10753v1 (CC BY 4.0)}
\begin{document}
\maketitle

\noindent\emph{Source and licence.} On a generalization of decomposable maps on C*-algebras. Version arXiv:2602.10753v1 (CC BY 4.0), distributed by arXiv under the Creative Commons Attribution 4.0 International licence, http://creativecommons.org/licenses/by/4.0/, as recorded in the arXiv licence metadata for this exact version and shown on the abstract page https://arxiv.org/abs/2602.10753v1. Full licence notice: \texttt{licenses/arxiv-2602.10753-CC-BY-4.0.txt}.\par\smallskip

\noindent\textit{Editorial scope.} Counted unit: the article's lemma lemma:Phi0ucp (source0.tex lines 510--516). The article's complete original proof of it is delivered with it (source0.tex lines 518--568, one \texttt{proof} environment of 51 lines). No mathematics is written, repaired or re-derived: the statement and the proof are the article's own lines, in the article's own notation, and no retained line is substituted or re-flowed.  No further source-local context is needed: every cross-reference of every retained line resolves inside the two delivered spans.  Nothing was omitted from any retained span, so the package carries no omission record.  No retained line contains a citation, so the wrapper carries no bibliography and no bibliography entry.  No retained line contains a figure environment, an image-inclusion command or drawing code, so no image is delivered and the package declares no asset dependency.  Editorial formatting only, in addition: an AMS \texttt{amsart} preamble that restates the article's own declarations and macros used by the retained lines, namely the environments \texttt{lemma} on separate counters, together with the standard packages those lines need.  The retained environments print in this document as Lemma 1 in order of appearance, whereas the article prints the same items with the numbering of its own full document.  The complete native source of the article as distributed in the arXiv e-print for this exact version is bundled unchanged as \texttt{sources/194425/source0.tex}.
\begin{lemma}\label{lemma:Phi0ucp}
    Let $\mathscr{A}$ be a nonunital C*-algebra, $\mathscr{M} \subset \mathscr{A}$ a C*-subalgebra, $H$ a Hilbert space and $T : \mathscr{M}\to B(H)$ a c.p.~map. Then, its extension $T^\sim : \mathscr{M} \oplus \complexes 1 \to B(H)$ given via
    \begin{equation}
        T^\sim ((a,z)) = T(a) + z \cbnorm{T} I,
    \end{equation}
    for $I$ the identity operator on $H$ and $\cbnorm{T}$ the CB norm of $T$, is completely positive.
\end{lemma}

\begin{proof}
    Denote $\mathscr{M}^\sim = \mathscr{M}\oplus\complexes 1$. We must show $T_{n}^{\sim} = \id{n}\otimes T^\sim$ is positive on $\matra{n}{\mathscr{M}^\sim}$, for all $n\in\naturals$. Let matrices $\left[x_{ij}\right]_{i,j} \in \matra{n}{\mathscr{M}}$ and $\Lambda \in \matr{n}$ be arbitrary. Then, any element $A \in \matra{n}{\mathscr{M}^\sim}$ admits a matrix representation
    \begin{equation}
        A = X + \Lambda \otimes \unit_\sim ,
    \end{equation}
    where $X = \left[(x_{ij},0)\right]_{i,j}$ and $\unit_\sim = (0,1)$ is a unit in $\mathscr{A}^\sim = \mathscr{A}\oplus \complexes 1$, the forced unitization of $\mathscr{A}$. Assume $A$ is a positive element in $\matra{n}{\mathscr{M}^\sim}$. Let $q : \mathscr{A}^\sim \to \complexes$ be a quotient map, $q ((a,z)) = z$. One checks $q$ is a *-homomorphism and a c.p.~map in consequence; hence $q_n = \id{n}\otimes q$ is positive. Since $q_n$ acts by restricting $A$ to its scalar part, $q_n (A) = \Lambda$, it must hold that $\Lambda$ is positive semidefinite. Define $\Lambda_\epsilon = \Lambda + \epsilon I_n$ and $A_\epsilon = X + \Lambda_\epsilon \otimes \unit_\sim$ for $\epsilon > 0$ and $I_n$ an identity matrix in $\matr{n}$. Since $A_\epsilon - A = \epsilon I_n \otimes \unit_\sim$ we see $A_\epsilon > A$ and so $\left(A_\epsilon\right)_\epsilon$ is a net of strictly positive elements, approximating $A$ in norm as $\epsilon \to 0$. Element $A_\epsilon$ is positive if and only if
    \begin{equation}
        (\Lambda_{\epsilon}^{-\frac{1}{2}}\otimes \unit_\sim)A_\epsilon (\Lambda_{\epsilon}^{-\frac{1}{2}}\otimes \unit_\sim) = Y_\epsilon + I_n \otimes \unit_\sim
    \end{equation}
    is positive as well, where we introduced
    \begin{equation}
        Y_\epsilon = (\Lambda_{\epsilon}^{-\frac{1}{2}}\otimes \unit_\sim)X (\Lambda_{\epsilon}^{-\frac{1}{2}}\otimes \unit_\sim).
    \end{equation}
    Positivity of $Y_\epsilon + I_n \otimes \unit_\sim$ is equivalent to $Y_\epsilon$ being self-adjoint and the spectrum $\spec{Y_\epsilon + I_n \otimes \unit_\sim}$ being nonnegative, i.e.
    \begin{equation}
        \spec{Y_\epsilon} \subset [-1,\infty ).
    \end{equation}
    Denote with $\spec{Y_\epsilon}_\pm$ the positive (nonnegative) and negative part of the spectrum, respectively. As $\mathscr{M}$ is a C*-algebra, it admits a well-defined functional calculus; the same is then true for $\mathscr{M}^\sim$ and $\matra{n}{\mathscr{M}^\sim}$. Let then $E$ be the spectral measure of $Y_\epsilon$. Define
    \begin{equation}
        Y_{\epsilon}^{\pm} = \int\limits_{\spec{Y_\epsilon}_\pm} |\lambda| \, dE(\lambda),
    \end{equation}
    so that $Y_{\epsilon}^{\pm}$ are positive, satisfy $Y_{\epsilon}^{+} Y_{\epsilon}^{-} = 0$ and $Y_\epsilon = Y_{\epsilon}^{+} - Y_{\epsilon}^{-}$. Since $Y_{\epsilon}^{-}$ is self-adjoint and its spectrum $\spec{Y_\epsilon}_-$ is contained in $[-1,0]$ we see $\| Y_\epsilon \| \leqslant 1$. It is an easy exercise to show that due to its definition, $Y_{\epsilon}^{\pm}$ have a structure $\left[(y_{ij}^{\pm}(\epsilon),0)\right]_{i,j}$ for some $\hat{y}_{\epsilon}^{\pm} = \left[y_{ij}^{\pm}(\epsilon)\right]_{i,j}\in \matra{n}{\mathscr{M}}$, so are contained in a subspace of $\matra{n}{\mathscr{M}^\sim}$ canonically isometrically isomorphic to $\matra{n}{\mathscr{M}}$; in particular $\left\|Y_{\epsilon}^{\pm}\right\| = \left\| \hat{y}_{\epsilon}^{\pm} \right\|$. Acting with $T_{n}^{\sim}$ we obtain
    \begin{equation}\label{eq:TnSimOnYplusUnit}
        T_{n}^{\sim}(Y_\epsilon + I_n \otimes \unit_\sim) = T_{n}^{\sim}(Y_{\epsilon}^{+}) - T_{n}^{\sim}(Y_{\epsilon}^{-}) + \cbnorm{T} I_n \otimes I,
    \end{equation}
    since $T_{n}^{\sim}(I_n \otimes \unit_\sim) = \left[ T^\sim ((0,1)) \right]_{i,j} = \cbnorm{T}\left[\delta_{ij}I\right]_{i,j}$ by definition of $T^\sim$. Notice
    \begin{equation}
        T_{n}^{\sim}(Y_{\epsilon}^{-}) = \left[ T^\sim ((y_{ij}^{-}(\epsilon),0)) \right]_{i,j} = \left[T(y_{ij}^{-}(\epsilon))\right]_{i,j} = T_n (\hat{y}^{-}_\epsilon)
    \end{equation}
    for $T_n = \id{n} \otimes T : \matra{n}{\mathscr{M}}\to \matra{n}{B(H)}$. Since we assumed $T$ was completely bounded and $\left\| Y_{\epsilon}^{-}\right\| \leqslant 1$, for any $n\in\naturals$ we have en estimation
    \begin{equation}
        \left\|T_{n}^{\sim}(Y_{\epsilon}^{-}) \right\| \leqslant \cbnorm{T} \left\|\hat{y}^{-}_{\epsilon}\right\| = \cbnorm{T} \left\| Y_{\epsilon}^{-} \right\| \leqslant \cbnorm{T}.
    \end{equation}
    Both elements $T_{n}^{\sim}(Y_{\epsilon}^{+})$, $T_{n}^{\sim}(Y_{\epsilon}^{-})$ are positive by the assumed complete positivity of $T$. In particular,
    \begin{equation}
        \spec{T_{n}^{\sim}(Y_{\epsilon}^{-})} \subset [0, \cbnorm{T}]
    \end{equation}
    which yields, with the use of the spectral theorem, that the spectrum
    \begin{equation}
        \spec{\cbnorm{T}I_n \otimes I - T_{n}^{\sim}(Y_{\epsilon}^{-})} = \{\cbnorm{T} - \lambda : \lambda \in [0, \cbnorm{T}] \}
    \end{equation}
    is contained in $[0,\cbnorm{T}]$, i.e.~$\cbnorm{T}I_n \otimes I - T_{n}^{\sim}(Y_{\epsilon}^{-})$ is positive. In consequence, the right hand side of \eqref{eq:TnSimOnYplusUnit} is positive for all $n$ and all $\epsilon > 0$. Denote $\mathcal{Y}_\epsilon = T_{n}^{\sim}(Y_\epsilon + I_n \otimes \unit_\sim)$. One checks with a direct computation that any map $\phi_n = \id{n}\otimes\phi$, where $\phi$ acts on $\mathscr{M}^\sim$, satisfies the intertwining relation
    \begin{equation}
        \phi_n \left((K \otimes \unit_\sim)X(L \otimes \unit_\sim)\right) = (K \otimes \unit_\sim) \phi_n(X) (L \otimes \unit_\sim)
    \end{equation}
    for any $K,L \in \matr{n}$ and $X \in \matra{n}{\mathscr{M}^\sim}$; this yields
    \begin{equation}
        T_{n}^{\sim} (A_\epsilon) = \left(\Lambda_{\epsilon}^{\frac{1}{2}}\otimes\unit_\sim\right) \mathcal{Y}_\epsilon \left(\Lambda_{\epsilon}^{\frac{1}{2}}\otimes\unit_\sim\right)
    \end{equation}
    defines a net of positive elements in $\matra{n}{B(H)}$. Net $\left(A_\epsilon\right)_\epsilon$ was norm convergent to $A$, so we have $\| A_\epsilon - A\| < \frac{\delta}{\|T^{\sim}_{n}\|}$ for some $\delta > 0$ and $\epsilon$ small enough; hence $\| T_{n}^{\sim} (A_\epsilon - A) \| < \delta$ and $\left(T_{n}^{\sim} (A_\epsilon)\right)_\epsilon$ converges to $T_{n}^{\sim}(A)$ in norm as $\epsilon \to 0$. Since $\matra{n}{B(H)}^+$ is norm closed, $T_{n}^{\sim}(A)$ must be positive, for any $A \in \matra{n}{\mathscr{M}^\sim}^+$ and any $n$. This shows $T^\sim$ is completely positive.
\end{proof}

\end{document}
